\chapter*{Abstract}
\addcontentsline{toc}{chapter}{Abstract}

Mutual information (MI) measures how much observing one variable reduces
uncertainty about another. This thesis considers whether MI between two
categorical variables differs across independent populations. A Normal
Wald test provides one approach, but small or sparse samples can distort
the estimated difference and its standard error, producing unreliable
false-positive rates and power.

This thesis develops an Expanded Welch--Satterthwaite test that retains
Wald's estimated difference and standard error and uses a Student \(t\)
distribution. Its degrees of freedom depend on how much the estimated
MI variance changes between samples. The thesis derives a first-order
approximation to that variability by accounting for changes in cell
probabilities, their margins and the resulting MI values.

Simulations check the moment approximations and compare false-positive
rates and power with Normal Wald and a sample-size degrees-of-freedom
rule. The mean and sampling-variance approximations have different errors,
and estimated degrees of freedom can remain well below the target based
on the actual sampling moments. An online-shopping dataset illustrates
interpretation on observed data.

Expanded Welch brings false-positive rates closer to the intended level
in selected small-sample populations, at a cost in nominal power. Its
power advantage after empirical calibration depends on which equal-MI
population supplies the cutoff and reverses in some comparisons when
that population changes. A larger cutoff also leaves sparse-table bias
and inaccurate standard errors unresolved. The results support selected
calibration improvements while identifying estimation errors that limit
reliable inference.
